\section{Hukum-hukum Dasar Aljabar Proposisi}

Sebagaimana rakitan disiplin ilmu matematika murni yang mendamba pemangkasan derivasi teorem secara instan menggunakan seperangkat aljabar linier, dunia percabangan *Boolean Logic* pun dibekali artileri senjata kembar. Kita menyepakatif daftar formula super "Jalan Pintas" yang sudah terlebih dahulu terbukti keabsahannya via puluhan tabel kebenaran di era purba. Himpunan senjata derivasi instan ini dikerucutkan dalam julukan agung \textbf{Hukum-hukum Aljabar Proposisi} \cite{rosen2019}.

Hukum-hukum Aljabar logika memecah diri ke dalam kelompok fungsional. Mari kita bedah gelombang demi gelombang senjata penyederhanaan ini.

\subsection{1. Hukum Pengaruh Absolut (Dominasi \& Identitas)}
Sebuah nilai raksasa Tautologi (T) atau Kontradiksi (F) bilamana disandingkan membajak variabel lemah $p$, seringkali berakhir mematikan sifat biner (Dominasi) atau malah sekedar tersenyum membias tanpa bayangan ikut campur (Identitas).

\begin{itemize}
    \item \textbf{Hukum Dominasi (Dominasi Mutlak / Null Laws)} \\
    Variabel selemah apapun bila bertatap muka Disjungsi ($Or$) dengan Tautologi Mutlak (T), akan ikut terpaksa ditarik bersinar terang menjadi T. \\ Sebaliknya mengikatkan konjungsi ($And$) kepada sang iblis kemustahilan false (F), akan tersedot jatuh ke jurang gelap F selamanya.
    \begin{align*}
        p \vee \mathbf{T} &\equiv \mathbf{T} \\
        p \wedge \mathbf{F} &\equiv \mathbf{F}
    \end{align*}

    \item \textbf{Hukum Identitas (Hukum Cermin Transparan)} \\
    Mengunggah koneksi AND dengan Tautologi (T) seolah cuma melenggang berkaca; eksistensi $p$ dipertahankan murni. Demikian pula bersinggungan relasi Disjungsi OR bermandikan sumur kosong (F).
    \begin{align*}
        p \wedge \mathbf{T} &\equiv p \\
        p \vee \mathbf{F} &\equiv p
    \end{align*}
\end{itemize}

\subsection{2. Hukum Konflik Internal (Negasi \& Komplemen)}
Ketika sesosok variabel terjebak perang saudara melawan *dopleganger* negasi palsu dari dirinya sendiri ($\neg p$).

\begin{itemize}
    \item \textbf{Hukum Negasi (Komplemen Akhir)} \\
    Menabrakkan diri bersama negasi sendiri memunculkan dua kasta absolut semesta (Tautologi utuh atau Kontradiksi kelam) seperti didiskusikan di Section 4.1.
    \begin{align*}
        p \vee \neg p &\equiv \mathbf{T} \quad \text{(Hukum Tiada Jalan Tengah)} \\
        p \wedge \neg p &\equiv \mathbf{F} \quad \text{(Hukum Mustahil Kontradiksi serentak)}
    \end{align*}

    \item \textbf{Hukum Involusi (Pembatalan Negasi Ganda)} \\
    Musuhnya si Musuh merangkul status persahabatan sejati kembali utuh.
    \begin{align*}
        \neg(\neg p) &\equiv p 
    \end{align*}
\end{itemize}

\subsection{3. Hukum Redundansi Berlebihan (Idempoten \& Absorpsi)}
Tujuan luhur aljabar logika adalah memutilasi kalimat yang berulang tumpang-tindih panjang berbusa.

\begin{itemize}
    \item \textbf{Hukum Idempoten (Sia-sia Kembar Identik)} \\
    Berkompromi panjang memutari variabel diri aslinya sendiri murni omong kosong pengulangan.
    \begin{align*}
        p \vee p &\equiv p \\
        p \wedge p &\equiv p
    \end{align*}

    \item \textbf{Hukum Absorpsi (Penyerapan Raksasa Hitam)} \\
    Aturan sakti paling elegan namun krusial dihafal di luar kepala. Bilamana variabel subjek merajut koneksi menyimpang dan tiba-tiba menemukan **anak struktur dirinya memantul ada di selongsong kurung persilangan**, si tubuh kurung beserta operator ganjil tetangganya dilumat habis ditelan subjek aslinya!
    \begin{align*}
        p \vee (p \wedge q) &\equiv p \\
        p \wedge (p \vee q) &\equiv p
    \end{align*}
\end{itemize}
Hafalkan hukum *Absorpsi* ini secara khusyuk. Percayalah, hukum inilah sang Juru Selamat algoritma pemendek sirkuit terhebat manakala dihidangkan soal ujian logika nan menyiksa mata \cite{wiki_proplogic}.
